% FILE: 04_implementation_surface.tex
% Project: livestream
% Thesis Role: real-time-process-broadcast-and-livestream-standards

\section{Implementation Surface}
\label{sec:livestream-impl}

\subsection{Source Files}
\label{sec:livestream-impl-sources}

All four broadcast modules are implemented in Node.js and reside in
\texttt{/Users/sac/process-intelligence/livestream/}. The autonomic controller and
browser visualizer reside in
\texttt{/Users/sac/process-intelligence/experiments/visualizer/}.

\textbf{aalst\_broadcaster.js.} Invokes \texttt{git log --pretty=format:"\%H|\%an|\%ai|\%s"}
via \texttt{execSync} in the \texttt{process-intelligence} repository directory.
Each output line is split on the pipe delimiter to extract hash, author, date, and
subject. The date string is wrapped in \texttt{new Date(\ldots).toISOString()} to
produce a UTC ISO 8601 timestamp. The resulting OCEL 2.0 event envelope is written
to stdout as a JSONL line. The module runs once and exits; it does not poll.

\textbf{master\_ocel\_stream.js.} Implements a recursive polling loop using
\texttt{setTimeout(poll, 1000)}: every second, it lists the replays directory, emits
any file not yet in the \texttt{seen} Set, and recurses. Emitted events are raw JSON
objects from the replay files, not yet re-enveloped in OCEL 2.0 format. The module
announces itself on startup with the banner
\texttt{--- AALST LIVE MASTER STREAM ACTIVE ---}.

\textbf{replay\_aggregator.js.} Reads all \texttt{.json} files in the replays
directory, extracts the nested path \texttt{input.ocel2.event\_log.events}, and
pushes all events into a master array. The result is written to
\texttt{/Users/sac/zoeapp/master\_conversation.ocel} as a JSON object with key
\texttt{ocel:events}. Event ID collisions (e.g., two replay files each containing
an event with ID \texttt{e1}) are handled by array concatenation without deduplication,
as noted in \texttt{conformance\_audit.md}: ``Collisions in event IDs (e1/e2) were
correctly handled by the aggregator.''

\textbf{agent\_conversation\_log.js.} Parses Gemini chatlog JSONL files from
\texttt{/Users/sac/.gemini/tmp/zoeapp/chats}. For each line, it checks for three
event types: \texttt{thoughts} entries are materialized as \texttt{Admissibility Guard}
intent events; \texttt{toolCalls} entries are materialized as typed activity events;
and \texttt{toolCalls[i].result} entries are materialized as \texttt{Event} type
records. The resulting array is sorted chronologically by timestamp and emitted as
a JSON object with key \texttt{ocel:events}. Note: the source file at line 18
contains a syntax error (missing opening brace: \texttt{\} catch(e) \}}), which means
the module cannot be executed in its current form without correction. This is an open
defect recorded in the PARTIAL status.

\textbf{autonomic.js} (\texttt{experiments/visualizer/}).
Exports a pure-JavaScript SHA-256 implementation (no external dependencies), the
\texttt{AutonomicController} class with \texttt{tick()} implementing the MAPE-K loop,
the \texttt{EWMACalculator} class (alpha\,=\,0.15, LCL\,=\,0.920), and the
\texttt{AuditLedger} class implementing the SHA-256 chain
$H_i = \text{SHA-256}(i \| t \| \text{case\_id} \| \text{payload} \| H_{i-1})$.
The file is governed under the Blue River Dam Epistemic Containment Protocol
(cited in the module header) and the Autonomic Knowledge Actuation specification.

\subsection{Commands}
\label{sec:livestream-impl-commands}

The broadcast modules are invoked directly via Node.js:
\begin{verbatim}
node /Users/sac/process-intelligence/livestream/aalst_broadcaster.js
node /Users/sac/process-intelligence/livestream/master_ocel_stream.js
node /Users/sac/process-intelligence/livestream/replay_aggregator.js
node /Users/sac/process-intelligence/livestream/agent_conversation_log.js
\end{verbatim}

The \texttt{aalst\_broadcaster.js} and \texttt{replay\_aggregator.js} modules are
one-shot; the \texttt{master\_ocel\_stream.js} module runs indefinitely until
interrupted. The \texttt{agent\_conversation\_log.js} module cannot execute in
its current form due to the syntax error on line 18.

The browser visualizer is opened by loading
\texttt{experiments/visualizer/index.html} in a browser. As of the Stream Director
audit (v30.1.2, 2026-06-01), the DOM binding mismatch between \texttt{index.html}
and \texttt{app.js} was declared remediated by the Stream Director agent, but this
remediation has not been independently re-audited.

\subsection{Data and Ontology Surfaces}
\label{sec:livestream-impl-data}

The OCEL 2.0 event envelope used by all four broadcast modules is defined by four
mandatory fields:
\begin{itemize}
  \item \texttt{ocel:activity} --- string naming the process activity;
  \item \texttt{ocel:timestamp} --- UTC ISO 8601 datetime string;
  \item \texttt{ocel:vmap} --- key-value map of event attributes (e.g., commit hash,
    author, message subject, tool call arguments);
  \item \texttt{ocel:omap} --- list of object identifiers associated with this event
    (e.g., the committing author, the tool name).
\end{itemize}

The full OCEL 2.0 property graph $\mathcal{L} = (\mathcal{E}, \mathcal{O},
\text{E2O}, \text{O2O}, \Delta)$ requires object entities ($\mathcal{O}$), E2O
links, and O2O links in addition to the event envelope. The Telemetry Auditor audit
(v30.1.2) confirmed that the browser visualizer telemetry emits zero object entities
and zero E2O/O2O links, achieving 0\% OCEL 2.0 compliance. The four JavaScript
broadcast modules partially satisfy the object requirement through \texttt{ocel:omap}
fields but do not construct a full object table.

\subsection{Templates and Rendered Artifacts}
\label{sec:livestream-impl-templates}

The primary rendered artifact is the \texttt{master\_conversation.ocel} file
materialized by \texttt{replay\_aggregator.js} at
\texttt{/Users/sac/zoeapp/master\_conversation.ocel}. This file served as the
input event log for the conformance audit reported in \texttt{conformance\_audit.md}
(Fitness\,=\,0.9880, Precision\,=\,1.0000 against the Blue River Dam thesis model).

The browser visualizer renders four interactive panels: a dynamic SVG Petri net
canvas, an A* alignment solver display, an LTL DECLARE constraint monitor, and a
real-time EWMA drift chart with LCL\,=\,0.920. The Cryptographic Process Ledger
panel renders block heights, hashes, and back-link pointers in a blockchain explorer
format. The M\&A Claims Verification Board evaluates EBITDA synergies
(\$5M target), compliance enforcement (\$3.5M target), and data room lineage
(\$1.2M target) in real time against the live stream, and rejects any claim whose
supporting ledger block has been tampered.
